\chapter{Secondary and Exploratory Results}
\label{app:secondary-results}

This appendix gives the complete results behind the secondary and exploratory
analyses summarized in Chapter~\ref{ch:results}. The same aggregation script
used for the main results assembles the tables. Point estimates come from
\texttt{measurements.csv}. The hierarchical confidence intervals in
Table~\ref{tab:all-slices} are imported from the record-level paired-bootstrap
analysis because they cannot be recomputed from the exported measurement table
alone.

As in the main Results chapter, a comparison is one reported combination of
prompting condition, entailment model, correctness rule, and pair of methods.
An AUROC difference is the first method's AUROC minus the second's, so a
positive value favors the first method.

\section{All semantic-versus-surface comparisons}
\label{app:all-comparisons}

Table~\ref{tab:all-slices} and Figure~\ref{fig:comparison-forest} report all 24
comparisons. Eleven hierarchical intervals exclude zero. Four favor semantic
entropy, and all four use the Qwen entailment model. Seven favor surface
entropy, and all seven use token-F1 correctness labels.

The two smaller Qwen-entailment-model differences that favor semantic entropy are below
the 0.025 variation observed across generation seeds. This is not a
contradiction. The seed range describes how much one method's AUROC changes
across seeds. A paired test compares two methods on the same generated answers,
so variation shared by both methods cancels when their scores are subtracted.
The seed range is therefore useful for judging one AUROC value, while the paired
interval is the correct check for a difference between methods.

% Generated by scripts/build_thesis_tables.py from docs/independent_analysis/data/measurements.csv. Do not edit by hand.
\begin{table}[htb]
\thisfloatsetup{capposition=above}
\centering
\footnotesize
\caption{All 24 comparisons: discrete semantic entropy minus surface entropy, with hierarchical intervals resampling both the nine units and the records within them. The final column names the direction where the interval excludes zero. Every comparison favoring semantic entropy uses the Qwen entailment model; every comparison favoring surface entropy uses token-F1 grading.}
\label{tab:all-slices}
\begingroup
\setlength{\tabcolsep}{3pt}
\begin{tabular}{@{}lllccl@{}}
\toprule
Condition & \makecell{Entailment\\model} & Grader & \makecell{AUROC\\difference} & \makecell{95\,\% CI\\(hierarchical)} & \makecell{Supported\\direction} \\
\midrule
\texttt{chat\_0shot} & \texttt{xlarge} & Qwen judge & +0.0204 & [-0.030, +0.059] & --- \\
\texttt{chat\_0shot} & \texttt{xlarge} & Llama judge & +0.0223 & [-0.033, +0.066] & --- \\
\texttt{chat\_0shot} & \texttt{xlarge} & token-F1 & -0.0506 & [-0.084, -0.022] & surface \\
\texttt{chat\_0shot} & \texttt{v3-large} & Qwen judge & +0.0147 & [-0.027, +0.050] & --- \\
\texttt{chat\_0shot} & \texttt{v3-large} & Llama judge & +0.0139 & [-0.039, +0.055] & --- \\
\texttt{chat\_0shot} & \texttt{v3-large} & token-F1 & -0.0587 & [-0.091, -0.029] & surface \\
\texttt{chat\_0shot} & \texttt{Qwen-72B} & Qwen judge & +0.0867 & [+0.060, +0.119] & SE \\
\texttt{chat\_0shot} & \texttt{Qwen-72B} & Llama judge & +0.0872 & [+0.056, +0.124] & SE \\
\texttt{chat\_0shot} & \texttt{Qwen-72B} & token-F1 & -0.0658 & [-0.144, -0.004] & surface \\
\texttt{default\_0shot} & \texttt{xlarge} & Qwen judge & +0.0032 & [-0.026, +0.028] & --- \\
\texttt{default\_0shot} & \texttt{xlarge} & Llama judge & +0.0084 & [-0.021, +0.033] & --- \\
\texttt{default\_0shot} & \texttt{xlarge} & token-F1 & -0.0263 & [-0.045, -0.012] & surface \\
\texttt{chat\_5shot} & \texttt{xlarge} & Qwen judge & -0.0020 & [-0.014, +0.008] & --- \\
\texttt{chat\_5shot} & \texttt{xlarge} & Llama judge & -0.0005 & [-0.011, +0.010] & --- \\
\texttt{chat\_5shot} & \texttt{xlarge} & token-F1 & -0.0174 & [-0.027, -0.007] & surface \\
\texttt{default\_5shot} & \texttt{xlarge} & Qwen judge & +0.0005 & [-0.007, +0.008] & --- \\
\texttt{default\_5shot} & \texttt{xlarge} & Llama judge & +0.0027 & [-0.006, +0.012] & --- \\
\texttt{default\_5shot} & \texttt{xlarge} & token-F1 & -0.0133 & [-0.022, -0.005] & surface \\
\texttt{default\_5shot} & \texttt{v3-large} & Qwen judge & +0.0014 & [-0.007, +0.009] & --- \\
\texttt{default\_5shot} & \texttt{v3-large} & Llama judge & +0.0047 & [-0.005, +0.015] & --- \\
\texttt{default\_5shot} & \texttt{v3-large} & token-F1 & -0.0130 & [-0.023, -0.004] & surface \\
\texttt{default\_5shot} & \texttt{Qwen-72B} & Qwen judge & +0.0119 & [+0.003, +0.021] & SE \\
\texttt{default\_5shot} & \texttt{Qwen-72B} & Llama judge & +0.0139 & [+0.003, +0.024] & SE \\
\texttt{default\_5shot} & \texttt{Qwen-72B} & token-F1 & -0.0064 & [-0.017, +0.004] & --- \\
\bottomrule
\end{tabular}
\endgroup
\end{table}


\begin{figure}[htb]
\centering
\includegraphics[width=\linewidth]{margin_forest.pdf}
\caption{All 24 semantic-versus-surface comparisons. The hierarchical
  intervals resample both the nine (dataset, model) units and the records within
  them. Filled markers identify the eleven intervals that exclude zero. The
  twelve comparisons that use a released cross-encoder with an LLM correctness
  judge appear at the bottom; none of their intervals excludes zero.}
\label{fig:comparison-forest}
\end{figure}

One token-F1 comparison does not exclude zero:
\texttt{default\_5shot} on the Qwen entailment model, with an interval of
$[-0.017, +0.004]$. This condition has the shortest answers and the strongest
entailment model, leaving little room for the two estimators to differ.

Every point estimate in Table~\ref{tab:all-slices} agrees to four decimal
places with the cell-level values behind Table~\ref{tab:sentence-length}. The
two tables are produced through different code paths, so this agreement checks
both paths.

\section{Secondary prompting-factorial analysis}
\label{sec:factorial-results}

Table~\ref{tab:factorial} reports the four prompting combinations on the
default entailment model, averaged over the two LLM correctness judges. The AUROC
difference between semantic entropy and surface entropy is $+0.021$ at 255.2
mean characters, $+0.006$ at 69.8, $-0.001$ at 11.4, and $+0.002$ at 11.2.
Every interval includes zero, and the largest difference is below the
seed-variation range.

% Generated by scripts/build_thesis_tables.py from docs/independent_analysis/data/measurements.csv. Do not edit by hand.
\begin{table}[htb]
\thisfloatsetup{capposition=above}
\centering
\footnotesize
\caption{The $2\times2$ on the canonical backend, mean of the two LLM judges. The isolating contrast is the pair of 0-shot corners, where demonstrations are held at zero and only the instruction changes. Every interval spans zero.}
\label{tab:factorial}
\begingroup
\setlength{\tabcolsep}{3pt}
\begin{tabular}{@{}lccrc@{}}
\toprule
Condition & Mean chars & Demos & SE $-$ surface & 95\,\% CI (units) \\
\midrule
\texttt{chat\_0shot} & 255.2 & 0 & +0.021 & [-0.029, +0.063] \\
\texttt{default\_0shot} & 69.8 & 0 & +0.006 & [-0.022, +0.031] \\
\texttt{chat\_5shot} & 11.4 & 5 & -0.001 & [-0.012, +0.008] \\
\texttt{default\_5shot} & 11.2 & 5 & +0.002 & [-0.006, +0.009] \\
\bottomrule
\end{tabular}
\endgroup
\end{table}


The two zero-shot conditions isolate the instruction because both omit worked
examples. Changing only the instruction changes mean answer length by a factor
of 3.7, while the AUROC difference falls from $+0.021$ to $+0.006$. In the two
five-shot conditions, changing only the instruction changes mean answer length
by 0.2 characters.

Two limits from Section~\ref{sec:conditions} make this a secondary analysis.
Both short-answer conditions are five-shot, so the design separates length from
example count only among the longer answers. The added conditions also exist
only on the default entailment model, so the factorial cannot test the entailment model with
the largest AUROC differences in Section~\ref{sec:clustering-sensitivity}.

\section{Exploratory long-form experiment}
\label{sec:longform-results}

This experiment is exploratory. It covers one biography task and one
correctness judge. The recorded results contain 500 biographies in each of nine
model--seed runs, but Section~\ref{sec:longform-design} explains that the exact entity list
cannot be recovered from the stored repository. The default entailment model was not
run, and this study uses a stricter paragraph-clustering rule than the original
study.

Table~\ref{tab:longform} reports those runs under claim-level grading by the
Qwen judge. Mean claim accuracy is 0.375.

\begin{table}[htb]
\thisfloatsetup{capposition=above}
\centering
\footnotesize
\caption{The long-form arm: nine cells, 500 biographies each, claim-level grading by the Qwen judge. Strict bidirectional entailment throughout, which is not the rule the original study applies to paragraphs (Section~\ref{sec:longform-design}).}
\label{tab:longform}
\begingroup
\setlength{\tabcolsep}{3pt}
\begin{tabular}{@{}lccc@{}}
\toprule
Backend & Discrete SE & Surface & Naive sample \\
\midrule
\texttt{exact-match} & 0.4965 & 0.4965 & 0.4973 \\
\texttt{nli-deberta-v3-base} & 0.5245 & 0.4965 & 0.4973 \\
\texttt{nli-deberta-v3-large} & 0.6703 & 0.4965 & 0.4973 \\
\textbf{\texttt{Qwen2.5-72B-Instruct}} & 0.7084 & 0.4965 & 0.4973 \\
\bottomrule
\end{tabular}
\endgroup
\par\vspace{2pt}\footnotesize Mean claim accuracy 0.375 over the nine cells. Surface and naive sample entropy are constant across backends because they never consult one; they are not performing badly but are structurally degenerate, since ten sampled biographies are never string-identical and every sample forms its own cluster. The canonical \texttt{deberta-v2-xlarge-mnli} was not applied to this arm.
\end{table}


Surface entropy scores 0.4965 and naive sample entropy 0.4973 under every
entailment model. Both are near chance. These baselines have a structural limit
for varied paragraph outputs: whenever all ten samples are different strings,
each forms its own group and the entropy reaches the same maximum value. The
aggregate AUROC values do not show that every record has a constant score.
In particular, a fully constant score would give AUROC exactly 0.5 with
average-rank handling of ties. These near-chance baselines provide only a weak
test of the semantic estimator.

The results also show that the choice of entailment model matters.
Discrete semantic entropy has an AUROC of 0.5245 with the base DeBERTa-v3 model,
0.6703 with the large variant, and 0.7084 with Qwen-72B.
This 0.184 range shows strong dependence on the entailment model,
but it is not a clean ranking of paragraph-entailment models. The experiment
uses \texttt{strict\_entailment} throughout, while the original paragraph
experiment uses the relaxed rule. A strict rule with a lower-recall model can
split answers too aggressively. The single Qwen correctness judge is also used
for clustering, so this experiment lacks a check with a different grader.
The original study reports its strongest paragraph result with DeBERTa, but
the different clustering rules mean that the current result is not a direct
test of that ranking.

\section{Secondary comparison of AUROC and AURAC}
\label{sec:aurac-disagreement}

AUROC is the primary metric in this thesis. Table~\ref{tab:aurac} provides a
secondary sensitivity analysis using the two rejection-accuracy definitions.
With AUROC, discrete semantic entropy leads surface entropy in 14 of 24
comparisons, with a mean difference of $+0.0016$. With the released
\texttt{aurac\_paper} definition, it leads in 8 of 24 and the mean difference
is $-0.0015$. With \texttt{aurac\_mean\_retained}, it leads in 6 of 24 and the
mean difference is $-0.0033$.

\input{includes/t_aurac.tex}

Across the 648 discrete-semantic-entropy cells at $T=1.0$,
\texttt{aurac\_paper} correlates with the generating model's accuracy at
$r = 0.971$, and \texttt{aurac\_mean\_retained} correlates at $r = 0.965$.
The corresponding correlation for AUROC is $r = 0.502$. The two AURAC measures
therefore reflect the model's accuracy much more strongly than AUROC does.
